\chapter{2018:Akshay Venkatesh}\label{ch2018akshayvenkatesh:akshay-venkatesh-fields-medal-2018-a-comprehensive-report-on-his-work-and-impact}

\begin{center}\rule{0.5\linewidth}{0.5pt}\end{center}

\textbf{Prepared:} July 2026

\textbf{Subject:} Akshay Venkatesh (b.\ 1981, New Delhi, India)

\textbf{Position:} Professor, School of Mathematics, Institute for Advanced Study, Princeton

\textbf{Doctoral Advisor:} Peter Sarnak (Princeton University)

\textbf{Fields Medal Citation:} \emph{``For his synthesis of analytic number theory, homogeneous dynamics, topology, and representation theory, which has resolved long-standing problems in areas such as the equidistribution of arithmetic objects.''}

\begin{center}\rule{0.5\linewidth}{0.5pt}\end{center}

\begin{prereqs}
Number theory (automorphic forms, L-functions, exponential sums), algebraic topology (cohomology of arithmetic groups, homological stability), basic representation theory (admissible representations of reductive groups over adeles), ergodic theory (mixing, equidistribution of flows on locally symmetric spaces).
\end{prereqs}

\begin{objectives}
Explain the subconvexity theorem of Michel--Venkatesh and its applications to equidistribution; describe the Cohen--Lenstra heuristics for function fields and their proof via topology; understand the connection between torsion in cohomology of arithmetic groups and analytic number theory.
\end{objectives}

\section{Problem History}

The equidistribution of arithmetic objects is one of the central themes connecting number theory, dynamics, and geometry. The classical question, dating to Gauss and Dirichlet, asks how integers are distributed among residue classes and how ideals of a number field are distributed among ideal classes. These questions generalize naturally to the distribution of periodic orbits on homogeneous spaces $\Gamma \backslash G$, where $\Gamma$ is an arithmetic lattice in a reductive group $G$. The fundamental challenge is to understand the limiting distribution of these objects as some natural parameter---such as the discriminant or the level---tends to infinity.

L-functions encode the arithmetic information governing these distributions. For an automorphic representation $\pi$ of $\operatorname{GL}_n$ over a number field $F$, the central value $L(1/2, \pi)$ controls the main term in equidistribution statements. The \textbf{subconvexity problem} asks for non-trivial bounds on this central value: $|L(1/2, \pi)| \ll C(\pi)^{1/4 - \delta}$ for some $\delta > 0$, where $C(\pi)$ is the analytic conductor.

The convexity bound follows from general Phragm\'en--Lindel\"of principles; any improvement has profound consequences for equidistribution, counting problems, and the spectral theory of automorphic forms. The subconvexity bound is the key technical ingredient needed to show that arithmetic sequences are well-distributed in families.

The subconvexity problem was known in special cases: Weyl (1921) for $\zeta(1/2+it)$, Burgess (1962) for GL(1) over finite fields, Duke--Friedlander--Iwaniec (1994) for GL(2) over $\mathbb{Q}$, and Michel--Venkatesh (2006) in a limited setting. But a uniform treatment over all number fields and all automorphic representations of GL(2) remained elusive. Previous approaches relied on special properties of the base field, preventing a unified treatment. The Duke--Friedlander--Iwaniec method used Hecke operators and Kloosterman sums in a way specific to $\mathbb{Q}$, while the analytic methods of Burgess worked only in the function field setting.

The \textbf{Lindel\"of hypothesis} for $L$-functions, which predicts the bound $|L(1/2+it, \pi)| \ll C(\pi)^\epsilon$ for any $\epsilon > 0$, is one of the most important unsolved problems in analytic number theory. Subconvexity is the first step toward the Lindel\"of hypothesis: any saving in the exponent beyond the convexity bound is called a subconvexity bound. The Lindel\"of hypothesis itself would imply, among other things, the strongest possible form of the prime number theorem and the best possible error term in the Bombieri--Vinogradov theorem.

The \textbf{Ramanujan conjecture} for $\operatorname{GL}(2)$ predicts that the Fourier coefficients of a weight $k$ modular form satisfy $|a_n| \le d(n) n^{(k-1)/2}$. The best known bound is due to Kim and Sarnak (2003), who achieved a saving of $1/7$ in the exponent. Subconvexity results are intimately related to the Ramanujan conjecture: each improvement in subconvexity typically yields an improvement in bounds on Fourier coefficients.

The \textbf{Kuznetsov trace formula} provides a fundamental tool for relating sums of Fourier coefficients of automorphic forms to sums involving Kloosterman sums. This formula, which generalizes the Petersson trace formula to non-compact quotients, is the technical foundation for many subconvexity results. Venkatesh's work extends the Kuznetsov trace formula approach to arbitrary number fields by using the theory of spherical functions and the spectral decomposition of automorphic forms on $\operatorname{GL}(2)$ over $\mathbb{A}_F$.

The \textbf{Cohen--Lenstra heuristics} (1983) predict the distribution of class groups of quadratic number fields, suggesting that class groups behave like random finite abelian groups weighted by $1/|\operatorname{Aut}|$. Proving any case of these heuristics required new methods connecting topology and arithmetic statistics. The difficulty lies in the fact that class groups are extremely difficult to control on average.

A related question concerns the growth of torsion in the cohomology of arithmetic groups. For a congruence subgroup $\Gamma$ of $\operatorname{SL}_2(\mathbb{Z})$, the first homology group $H_1(\Gamma, \mathbb{Z})$ is a finitely generated abelian group whose torsion subgroup grows exponentially with the level of $\Gamma$. The precise rate of growth and its connection to $L$-functions were not understood.

The interplay between these different threads---subconvexity, equidistribution, class group statistics, and torsion in cohomology---forms the backdrop to Venkatesh's work. Each problem is deep in its own right, but Venkatesh's genius was to see that they are connected by common geometric and dynamical principles. The synthesis of these ideas led to breakthroughs that none of the individual approaches could have achieved alone.

An important precursor to Venkatesh's work was the subconvexity result of Michel and Venkatesh (2006), which handled the case of $L$-functions of quadratic twist families over $\mathbb{Q}$. This earlier result used a trace formula approach that was powerful but limited in scope. The key difficulty in extending to arbitrary number fields was the lack of a sufficiently flexible equidistribution result for the relevant orbits on the adelic quotient. Venkatesh's subsequent work addressed this difficulty head-on by developing new tools from homogeneous dynamics that could be applied uniformly across all base fields.

The \textbf{Ramanujan conjecture} for $\operatorname{GL}(2)$, while not needed for the subconvexity result, provides important context. The conjecture predicts that the Fourier coefficients of a weight $k$ modular form satisfy $|a_n| \le d(n) n^{(k-1)/2}$. The best known bound is due to Kim and Sarnak (2003), who achieved an saving of $1/7$ in the exponent. Subconvexity results are intimately related to the Ramanujan conjecture: each improvement in subconvexity typically yields an improvement in bounds on Fourier coefficients, and vice versa.

\section{How Was It Solved}

Venkatesh's approach, developed in several stages with different collaborators, combined analytic number theory, homogeneous dynamics, and algebraic topology in unexpected ways. The key insight was that problems in number theory could be approached through the geometric and dynamical properties of homogeneous spaces.

\textbf{The Michel--Venkatesh subconvexity theorem.} For any automorphic representation $\pi$ of $\operatorname{GL}(2)$ over a number field $F$, the $L$-function satisfies $|L(1/2, \pi)| \ll C(\pi)^{1/4 - \delta}$ with $\delta > 0$ depending only on $[F:\mathbb{Q}]$. The proof proceeds by expressing $L(1/2, \pi)$ as a period integral of an automorphic form on $\operatorname{GL}(2) \times \operatorname{GL}(1)$, then bounding this integral using equidistribution of certain flows on the quotient.

The amplification method involves introducing a carefully chosen sum over auxiliary automorphic forms. By summing $L(1/2, \pi \otimes \chi)$ over characters $\chi$ with explicit weights, the main term is amplified while the error terms remain controlled. This allows one to extract a bound for individual $L$-functions from the collective behavior of the family.

The kernel function construction uses the theory of spherical functions on symmetric spaces. For each automorphic representation $\pi$, one constructs a test function on the adelic group whose Fourier transform captures the central value $L(1/2, \pi)$. This connects the analytic problem to the geometry of the quotient $\Gamma \backslash G$ and allows the use of dynamical methods.

\textbf{Equidistribution of torus orbits.} With Einsiedler, Lindenstrauss, and Michel, Venkatesh proved that periodic torus orbits in $\operatorname{SL}_3(\mathbb{Z}) \backslash \operatorname{SL}_3(\mathbb{R})$ become equidistributed as the discriminant of the associated torus grows. This was the first equidistribution result for a non-compact homogeneous space of higher rank. The proof combines the ergodic theory of flows on homogeneous spaces with arithmetic estimates on the distribution of tori.

The significance of this result lies in its connection to the distribution of rational points on homogeneous varieties. A torus $T$ over $\mathbb{Q}$ embeds into $\operatorname{GL}_n$, and the periodic orbit corresponding to $T$ represents the set of rational points on the associated homogeneous variety. The equidistribution theorem shows that these rational points become uniformly distributed as the torus varies.

The proof uses a \textbf{mixing estimate} for the geodesic flow on $\operatorname{SL}_3(\mathbb{Z}) \backslash \operatorname{SL}_3(\mathbb{R})$. Unlike the compact case, where mixing is immediate from the spectral gap, the non-compact setting requires careful analysis of the continuous spectrum. The key technical innovation is a \textbf{quantitative effective mixing estimate} that controls the rate of equidistribution in terms of the discriminant of the torus.

The connection to number theory comes through the \textbf{torus period formula}: the period of a Hecke eigenform $f$ along a torus orbit $\mathcal{O}_T$ equals $L(1/2, f \otimes \chi_T)$, where $\chi_T$ is the Hecke character associated to $T$. This identity, due to Jacquet and Shalika, converts the analytic problem of bounding $L$-values into the geometric problem of studying equidistribution.

\textbf{The Cohen--Lenstra heuristics for function fields.} With Ellenberg and Westerland, Venkatesh proved that the Cohen--Lenstra heuristics hold for class groups of function fields over finite fields. The proof uses the stable homology of moduli spaces of curves, combined with Harer's stability theorem and the Madsen--Weiss theorem. The key insight is that the arithmetic statistics problem translates into a problem about the topology of these moduli spaces. Specifically, the Cohen--Lenstra weight $1/|\operatorname{Aut}(A)|$ arises naturally from the cohomology of the symmetric group acting on the moduli space.

The proof proceeds in several steps. First, one counts the number of function fields with a given class group by counting suitable covers of a fixed curve. The number of such covers is governed by the homology of the mapping class group. Second, one applies Harer's stability theorem to identify the stable homology of the mapping class group. Third, one uses the Madsen--Weiss theorem to compute this stable homology in terms of characteristic classes. The final step is to show that the resulting counting formula matches the Cohen--Lenstra prediction.

The proof also uses the theory of \textbf{random matrices} and the statistics of eigenvalues of random matrices over finite groups. The Cohen--Lenstra heuristic itself was motivated by numerical evidence and by analogies with random matrix theory. The fact that the distribution of class groups matches the distribution of cokernels of random matrices over $\mathbb{Z}$ is a remarkable coincidence that is explained by the topological methods of Venkatesh, Ellenberg, and Westerland.

The connection to the Madsen--Weiss theorem is particularly beautiful. This theorem, which is one of the deepest results in the topology of moduli spaces, identifies the stable homology of the mapping class group with the cohomology of an infinite loop space. The fact that this purely topological result has arithmetic consequences is one of the most striking examples of the unity of mathematics.

The proof of the Madsen--Weiss theorem itself, due to Galatius, Madsen, Tillmann, and Weiss, uses techniques from higher algebraic topology, including the theory of $E_\infty$-ring spectra and the cobordism hypothesis. The application to the Cohen--Lenstra heuristics provides a concrete arithmetic motivation for these abstract topological constructions, demonstrating the deep interplay between different areas of mathematics.

An important aspect of the proof is the use of \textbf{harmonic analysis on finite groups}. The Cohen--Lenstra weight $1/|\operatorname{Aut}(A)|$ can be expressed as an integral over the group of characters of $A$. This integral representation allows one to use the tools of harmonic analysis to study the distribution of class groups, and it connects the problem to the theory of random walks on finite groups.

The techniques developed by Venkatesh, Ellenberg, and Westerland have been applied to other problems in arithmetic statistics, including the study of $S$-integer points on algebraic groups, the distribution of ranks of quadratic twists of elliptic curves, and the statistics of Selmer groups. These applications demonstrate the versatility of the topological approach to arithmetic statistics.

\section{Impact of This Discovery}

Venkatesh's work transformed analytic number theory by demonstrating that homogeneous dynamics is a natural and powerful tool for subconvexity. The Michel--Venkatesh theorem is now a standard reference, and subsequent work has extended these techniques to higher-rank groups and general number fields.

The Cohen--Lenstra result opened a new paradigm in arithmetic statistics: using topological methods to study the distribution of arithmetic objects. This approach has inspired dozens of subsequent papers, including extensions to cubic and higher-degree extensions of number fields.

The effective equidistribution results developed with Einsiedler and Margulis have applications beyond number theory, including quantum chaos, ergodic theory, and the geometry of locally symmetric spaces.

The connection between number theory and dynamics established by this work continues to yield new results. The methods have been applied to problems in additive combinatorics, thin groups, and the geometry of numbers. The subconvexity bound has become a crucial ingredient in counting problems, such as counting rational points of bounded height on homogeneous varieties.

The Cohen--Lenstra program has been extended to many settings beyond function fields. Bhargava, Shankar, and their collaborators have made progress on average bounds for class groups of number fields using the geometry of numbers and the representation theory of symmetric groups. The statistical behavior of class groups remains one of the most active areas in arithmetic statistics. The methods have also been applied to study the distribution of Selmer groups of elliptic curves, providing new evidence for the Birch and Swinnerton-Dyer conjecture.

The work of Venkatesh has also influenced the theory of \textbf{automorphic forms} and their $L$-functions. The subconvexity bound has been used to prove equidistribution results for Hecke eigenvalues, which have applications to the spectral theory of automorphic forms and to the theory of quantum unique ergodicity. The methods have been extended to handle $L$-functions of higher-degree automorphic representations, opening new avenues for research.

In the theory of \textbf{arithmetic geometry}, the Cohen--Lenstra heuristics have been used to predict the distribution of $S$-class groups, Mordell--Weil groups, and other arithmetic invariants. The topological methods developed by Venkatesh and his collaborators provide a systematic framework for studying these distributions, and the results have applications to the conjecture of Goldfeld on the average size of $S$-Selmer groups of quadratic twists of elliptic curves.

\section{Modern Views and Improvements}

Venkatesh continues to develop his conjectures on beyond endoscopy, a program proposed by Langlands that aims to prove functoriality of automorphic representations using $L$-functions and trace formulas. The limit multiplicity formula and the spectral decomposition of $L$-functions are active research directions. The beyond endoscopy program envisions using the subconvexity bound as a starting point for proving the general functoriality conjecture.

The connection between torsion in cohomology and analytic number theory is being pursued by many researchers. Recent work has revealed deep connections to the behavior of $L$-functions in the critical strip and to the theory of Iwasawa modules. The growth of torsion in homology provides new invariants for studying arithmetic groups, and the precise relationship between torsion and $L$-values continues to be an area of active investigation.

Higher-rank analogues of the equidistribution results remain a frontier. Extending the theorem beyond $\operatorname{GL}(2)$ requires new understanding of higher-rank homogeneous spaces and the representation theory of more general reductive groups. The interactions with the theory of automorphic representations continue to yield new results in this rapidly developing area of mathematics.

The development of \textbf{perfectoid spaces} by Scholze has opened new avenues for studying $L$-functions in the $p$-adic setting, complementing Venkatesh's work in the archimedean setting. The combination of these techniques promises to lead to a deeper understanding of the analytic behavior of $L$-functions and their arithmetic applications.

Recent work by Shin and Templier on the equidistribution of $\operatorname{GL}(n)$ Hecke orbits in the Sato--Tate theorem uses techniques inspired by Venkatesh's approach. The Sato--Tate conjecture, now a theorem for $\operatorname{GL}(2)$ over totally real fields, describes the distribution of angles of Frobenius elements in families of modular forms. The proof uses the subconvexity bound as a key ingredient.

The development of \textbf{condensed mathematics} by Clausen and Scholze has provided new foundations for functional analysis and algebraic topology that may eventually have applications to the study of automorphic forms. While the connections to Venkatesh's work are still being explored, the potential for new applications is significant.

The interaction between Venkatesh's work and the theory of \textbf{random matrix theory} continues to be a fruitful area of research. The distribution of zeros of $L$-functions is conjectured to follow the statistics of eigenvalues of random matrices from classical compact groups, and the subconvexity bound provides evidence for this conjecture. The methods of Venkatesh have been used to prove results about the distribution of low-lying zeros of $L$-functions, providing new evidence for the Katz--Sarnak philosophy.

\begin{center}\rule{0.5\linewidth}{0.5pt}\end{center}

\section*{Summary}
\addcontentsline{toc}{section}{Summary}

Venkatesh synthesized analytic number theory, homogeneous dynamics, and topology to prove subconvexity bounds for $\operatorname{GL}(2)$ $L$-functions, equidistribution of torus orbits, and the Cohen--Lenstra heuristics for function fields. His work demonstrates the deep unity of mathematics and has opened new directions in arithmetic statistics, automorphic forms, and the geometry of numbers.
